\documentclass[10pt,a4paper]{amsart}
\usepackage[margin=19mm]{geometry}
\usepackage{amsmath,amssymb,mathtools}
\usepackage[scr=rsfs,cal=euler]{mathalfa}
\usepackage{xfrac}
\usepackage[unicode,hidelinks,bookmarks=false]{hyperref}
\newcommand{\ie}{i.e.\ }
\newcommand{\cf}{cf.\ }
\newcommand{\resp}{respectively\ }
\newcommand{\Subsection}{\subsection}
\providecommand{\nobreakdash}{\nobreak\hbox{-}}
\setlength{\emergencystretch}{2em}
\multlinegap0pt
\usepackage[shortlabels]{enumitem}
\usepackage{amssymb}
\usepackage{mathtools}
\usepackage{tikz}
\usepackage{algorithm}\usepackage{algpseudocode}
\usepackage{float}
\newtheorem{theorem}{Theorem}
\newcommand{\calC}{\mathcal{C}}
\newcommand{\calL}{\mathcal{L}}
\newcommand{\calT}{\mathcal{T}}
\newcommand{\one}{\mathbf{1}}
\title{Block Structure and Spectrum of Zero-Divisor Graphs of Lipschitz Quaternion Rings Modulo $n$}
\author{Bilal Ahmad Rather}
\date{Source: arXiv:2603.20947v2 (CC BY 4.0)}
\begin{document}
\maketitle

\noindent\emph{Source and licence.} Block Structure and Spectrum of Zero-Divisor Graphs of Lipschitz Quaternion Rings Modulo $n$. Version arXiv:2603.20947v2 (CC BY 4.0), distributed by arXiv under the Creative Commons Attribution 4.0 International licence, http://creativecommons.org/licenses/by/4.0/, as recorded in the arXiv licence metadata for this exact version and shown on the abstract page https://arxiv.org/abs/2603.20947v2. Full licence notice: \texttt{licenses/arxiv-2603.20947-CC-BY-4.0.txt}.\par\smallskip

\noindent\textit{Editorial scope.} Counted unit: the article's theorem thm:fullreducedspectrum (source0.tex lines 532--568). The article's complete original proof of it is delivered with it (source0.tex lines 570--645, one \texttt{proof} environment of 76 lines). No mathematics is written, repaired or re-derived: the statement and the proof are the article's own lines, in the article's own notation, and no retained line is substituted or re-flowed.  Retained uncounted as the source-local context the proof needs: lines 470--489.  Nothing was omitted from any retained span, so the package carries no omission record.  No retained line contains a citation, so the wrapper carries no bibliography and no bibliography entry.  No retained line contains a figure environment, an image-inclusion command or drawing code, so no image is delivered and the package declares no asset dependency.  Editorial formatting only, in addition: an AMS \texttt{amsart} preamble that restates the article's own declarations and macros used by the retained lines, namely the environments \texttt{theorem} on separate counters, together with the standard packages those lines need.  The retained environments print in this document as Theorem 1 in order of appearance, whereas the article prints the same items with the numbering of its own full document.  The complete native source of the article as distributed in the arXiv e-print for this exact version is bundled unchanged as \texttt{sources/196897/source0.tex}.
	\begin{theorem}\label{thm:charfactor}
		Let $p$ be an odd prime, and let $A_p=A(G_p)$. For each type class $\calC_{L,M}$, let
		$\one_{L,M}\in \mathbb{R}^{V(G_p)}$ be its indicator vector. Let $U=\operatorname{span}\{\one_{L,M}:\ (L,M)\in \calT\},$ be the class-constant subspace, and for each class $\calC_{L,M}$ let
		$$
		W_{L,M}=\left\{x\in \mathbb{R}^{V(G_p)}:\ \operatorname{supp}(x)\subseteq \calC_{L,M},\ \sum_{v\in \calC_{L,M}}x_v=0\right\}.
		$$
		Then the following hold.
		\begin{enumerate}[label=\textup{(\roman*)},noitemsep]
			\item $\mathbb{R}^{V(G_p)}=U\oplus\bigoplus_{(L,M)\in\calT}W_{L,M}$ is an $A_p$-invariant direct sum decomposition.
			\item On $U$, the matrix of $A_p$ relative to the basis $\{\one_{L,M}\}_{(L,M)\in\calT}$ is $B_p:=(p-1)H_p-D_p.$ 
			\item If $L\neq M$, then $A_p|_{W_{L,M}}=0$.
			\item If $L=M$, then $A_p|_{W_{L,L}}=-I$.
		\end{enumerate}
		Consequently, the characteristic polynomial of $A_p$ is
		$$
		\chi_{A_p}(\lambda)
		=
		\lambda^{\,p(p+1)(p-2)}(\lambda+1)^{(p+1)(p-2)}\chi_{B_p}(\lambda).
		$$
	\end{theorem}

	\begin{theorem}\label{thm:fullreducedspectrum}
		Let $p$ be an odd prime and Let
		$
		\Delta_p=4p^4-4p^2-8p+9,~
		\Theta_p=p^4-8p+8.
		$
		Define
		$$
		\lambda_{\pm}=p^2-p-\frac12\pm\frac12\sqrt{\Delta_p},
		\quad
		\text{and}
		\quad
		\mu_{\pm}=\frac{p(p-2)\pm\sqrt{\Theta_p}}{2}.
		$$
		Then the spectrum of the reduced matrix $B_p=(p-1)H_p-D_p$ is
		$$
		\left\{
		\lambda_+,\lambda_-,
		\mu_+^{[p]},\mu_-^{[p]},
		\bigl(-p(p-1)\bigr)^{[p]},
		(p-1)^{[p(p-1)/2]},
		\bigl(-(p-1)\bigr)^{[(p+1)(p-2)/2]}
		\right\},
		$$
		where the superscript $[m]$ denotes multiplicity $m$. Equivalently,
		$$
		\begin{aligned}
			\chi_{B_p}(\lambda)
			={}&(\lambda-\lambda_+)(\lambda-\lambda_-)
			(\lambda-\mu_+)^p(\lambda-\mu_-)^p
			(\lambda+p(p-1))^p\\
			&\times(\lambda-(p-1))^{p(p-1)/2}
			(\lambda+(p-1))^{(p+1)(p-2)/2}.
		\end{aligned}
		$$
		Together with Theorem~\ref{thm:charfactor}, this gives the complete adjacency spectrum of $G_p$.
	\end{theorem}

	\begin{proof}
		Write $q=p+1$ and $r=p-1$. Label the lines of $\calL$ by $1,\ldots,q$ and identify $\mathbb R^{\calT}$ with the space of real $q\times q$ matrices $X=(x_{ij})$. If
		$$
		R_j=\sum_{\ell=1}^q x_{j\ell},
		\qquad
		C_i=\sum_{k=1}^q x_{ki},
		$$
		then the definition of $H_p$ gives
		$$
		(H_pX)_{ij}=R_j+C_i-x_{ji}.
		$$
		The subtraction of $x_{ji}$ removes the unique type $(j,i)$ counted in both incidence sums. Since $D_p$ keeps only diagonal coordinates,
		$$
		(B_pX)_{ij}=r(R_j+C_i-x_{ji})-\delta_{ij}x_{ii}.
		ag{4.1}
		$$
		
		Let $\one\in\mathbb R^q$ be the all-ones vector. We first consider
		$$
		\mathscr A=\{X:X^{\top}=-X,\ X\one=0\}.
		$$
		For $X\in\mathscr A$, both row and column sums vanish and $x_{ji}=-x_{ij}$, so (4.1) gives $B_pX=rX$. The dimension is
		$$
		\dim\mathscr A=\frac{(q-1)(q-2)}{2}=\frac{p(p-1)}{2}.
		$$
		Next let
		$$
		\mathscr S=\{X:X^{\top}=X,\ \operatorname{diag}(X)=0,\ X\one=0\}.
		$$
		Then (4.1) gives $B_pX=-rX$. The space of symmetric zero-diagonal matrices has dimension $q(q-1)/2$, and the row-sum map has rank $q$ for $q\ge3$; hence
		$$
		\dim\mathscr S=\frac{q(q-3)}{2}=\frac{(p+1)(p-2)}{2}.
		$$
		
		The two-dimensional space of matrices having one constant value on the diagonal and another off the diagonal is invariant. Relative to these two values, (4.1) acts as
		$$
		Q_p=
		\begin{pmatrix}
			p-2&2p(p-1)\\
			2(p-1)&(2p-1)(p-1)
		\end{pmatrix},
		$$
		whose eigenvalues are $\lambda_+$ and $\lambda_-$.
		
		It remains to account for the standard-coordinate part. Choose $u\in\one^{\perp}$ and consider
		$$
		\mathscr V(u)=\operatorname{span}\{u\one^{\top},\ \one u^{\top},\ \operatorname{diag}(u)\}.
		$$
		Equation (4.1) shows that this space is invariant. Relative to the displayed ordered basis, the representing matrix is
		$$
		M_p=
		\begin{pmatrix}
			0&p(p-1)&p-1\\
			p(p-1)&0&p-1\\
			-1&-1&-p
		\end{pmatrix}.
		$$
		The vector $(1,-1,0)^{\top}$ is an eigenvector with eigenvalue $-p(p-1)$. On the invariant plane where the first two coordinates agree, $M_p$ reduces to
		$$
		\begin{pmatrix}
			p(p-1)&p-1\\
			-2&-p
		\end{pmatrix},
		$$
		whose characteristic polynomial is
		$$
		\lambda^2-p(p-2)\lambda-(p-1)(p^2-2).
		$$
		Its two roots are exactly $\mu_+$ and $\mu_-$. Taking an orthogonal basis $u_1,\ldots,u_p$ of $\one^{\perp}$ therefore supplies each of these three eigenvalues with multiplicity $p$.
		
		The spaces just described are mutually orthogonal for the Frobenius inner product after the choice of an orthogonal basis of $\one^{\perp}$. Their dimensions add to
		$$
		2+\frac{p(p-1)}{2}+\frac{(p+1)(p-2)}{2}+3p=(p+1)^2.
		$$
		Thus they exhaust the reduced type space. Collecting the eigenvalues and their multiplicities gives the stated spectrum and characteristic polynomial.
	\end{proof}

\end{document}
